\documentclass[10pt,a4paper]{amsart}
\usepackage[margin=19mm]{geometry}
\usepackage{amsmath,amssymb,mathtools}
\usepackage[scr=rsfs,cal=euler]{mathalfa}
\usepackage{xfrac}
\usepackage[unicode,hidelinks,bookmarks=false]{hyperref}
\newcommand{\ie}{i.e.\ }
\newcommand{\cf}{cf.\ }
\newcommand{\resp}{respectively\ }
\newcommand{\Subsection}{\subsection}
\providecommand{\nobreakdash}{\nobreak\hbox{-}}
\setlength{\emergencystretch}{2em}
\multlinegap0pt
\usepackage{amssymb}
\usepackage{xcolor}
\newtheorem{lemma}{Lemma}
\newtheorem{theorem}{Theorem}
\def\r{\mathbb R}
\DeclareMathOperator{\sech}{sech}
\def\v{\mathcal V}
\title{Curves in Riemannian Manifolds Making Prescribed Angles With Torse-Forming Vector Fields}
\author{Muhittin Evren Aydin, Esra Dilmen, Busra Karakaya}
\date{Source: arXiv:2603.28072v1 (CC BY 4.0)}
\begin{document}
\maketitle

\noindent\emph{Source and licence.} Curves in Riemannian Manifolds Making Prescribed Angles With Torse-Forming Vector Fields. Version arXiv:2603.28072v1 (CC BY 4.0), distributed by arXiv under the Creative Commons Attribution 4.0 International licence, http://creativecommons.org/licenses/by/4.0/, as recorded in the arXiv licence metadata for this exact version and shown on the abstract page https://arxiv.org/abs/2603.28072v1. Full licence notice: \texttt{licenses/arxiv-2603.28072-CC-BY-4.0.txt}.\par\smallskip

\noindent\textit{Editorial scope.} Counted unit: the article's theorem thm42 (source0.tex lines 327--333). The article's complete original proof of it is delivered with it (source0.tex lines 334--352, one \texttt{proof} environment of 19 lines). No mathematics is written, repaired or re-derived: the statement and the proof are the article's own lines, in the article's own notation, and no retained line is substituted or re-flowed.  Retained uncounted as the source-local context the proof needs: lines 240--242; lines 245--247; lines 249--261; lines 263--291; lines 295--297; lines 308--321; lines 309--311.  Nothing was omitted from any retained span, so the package carries no omission record.  No retained line contains a citation, so the wrapper carries no bibliography and no bibliography entry.  No retained line contains a figure environment, an image-inclusion command or drawing code, so no image is delivered and the package declares no asset dependency.  Editorial formatting only, in addition: an AMS \texttt{amsart} preamble that restates the article's own declarations and macros used by the retained lines, namely the environments \texttt{lemma}, \texttt{theorem} on separate counters, together with the standard packages those lines need.  The retained environments print in this document as Lemma 1, Theorem 1 in order of appearance, whereas the article prints the same items with the numbering of its own full document.  The complete native source of the article as distributed in the arXiv e-print for this exact version is bundled unchanged as \texttt{sources/197406/source0.tex}.
\begin{equation}
f\langle \v, T\rangle+\omega(T)=0. \label{cdimm}
\end{equation}

\begin{lemma}\label{consleng}
Let $\gamma $ be a Frenet curve in a Riemannian manifold and $\v$ a unit torse-forming vector field along $\gamma$. Its potential function vanishes along $\gamma$ if and only if $\v$ is parallel along $\gamma$.
\end{lemma}

Assume now that $\gamma$ is a PA curve in $M^m$ with $(\v,\theta)$, where $\v$ is a torse-forming vector field. By its definition, there are smooth functions $\lambda_i:\,I\to \r$, $0\leq i \leq m-2$, defined along $\gamma$ such that
\begin{equation}
\left. 
\begin{array}{l}
\v=\lambda_0 T+\cos\theta N+\lambda_1 B_1+\ldots \lambda_{m-2}B_{m-2}, \\ 
1=\cos^2\theta+\sum_{i=0}^{m-2}\lambda_i^2.
\end{array}%
\right. \label{trs41}
\end{equation}
With this decomposition of $\v$, equation \eqref{cdimm} becomes
\begin{equation}
f\lambda_0+\omega(T)=0. \label{cdimm2}
\end{equation}

Differentiating $\v$ along $\gamma$ in the first expresion of \eqref{trs41} and using the Frenet formulas yields
\begin{equation}
\left. 
\begin{array}{l}
\nabla_T\v= (\lambda'_0-\cos\theta \kappa_1)T\\
+((\cos \theta)'+\lambda_0\kappa_1-\lambda_1\kappa_2)N \\
+(\lambda_1'+\cos\theta\kappa_2-\lambda_2\kappa_3)B_1\\
+(\lambda_i'+\lambda_{i-1}\kappa_{i+1}-\lambda_{i+1}\kappa_{i+2})B_i, \quad (2\leq i \leq m-3),\\
+(\lambda'_{m-2}+\lambda_{m-3}\kappa_{m-1})B_{m-2} .
\end{array}%
\right. \label{trs42}
\end{equation}
On the other hand, since $\v$ is a torse-forming vector field along $\gamma$, we have
$$
\nabla_T\v=(f+\lambda_0\omega(T))T+\omega(T)\cos\theta N+\omega(T)\sum_{i=1}^{m-2}\lambda_iB_{i}.
$$
Comparing this expression with \eqref{trs42} and then using \eqref{cdimm2}, we obtain
\begin{equation}
\left. 
\begin{array}{r}
f(1-\lambda_0^2)-\lambda'_0+ \kappa_1\cos\theta=0\\
-f\lambda_0\cos\theta+\theta'\sin \theta-\lambda_0\kappa_1+\lambda_1\kappa_2=0 \\
-f\lambda_0\lambda_1-\lambda_1'+\lambda_2\kappa_3-\kappa_2\cos\theta=0\\
-f\lambda_0\lambda_i-\lambda_i'-\lambda_{i-1}\kappa_{i+1}+\lambda_{i+1}\kappa_{i+2}=0 \\
-f\lambda_0\lambda_{m-2}-\lambda'_{m-2}-\lambda_{m-3}\kappa_{m-1}=0,
\end{array}%
\right. \label{trs43}
\end{equation}
where $2\leq i \leq m-3$.

\begin{theorem}\label{prop41}
Let $\gamma $ be a Frenet curve in a connected orientable Riemannian $m$-manifold. Then, the curve $\gamma$ is a PA curve with $(\v,\theta)$, where $\v$ is a torse-forming vector field with potential function $f$ and generating form $\omega$, if and only if there are smooth functions $\theta,\lambda_0,\ldots\lambda_{m-2}$ along $\gamma$ such that $\v$ is given by \eqref{trs41}, the curvatures of $\gamma$ satisfy the system \eqref{trs43} and $\omega(T)=-f\langle \v,T\rangle$, where $T$ is the unit tangent vector.
\end{theorem}

In the particular $3$-dimensional case, writing $B_1=B$, $\kappa_1=\kappa$ and $\kappa_2=\tau$, the decomposition of $\v$ in \eqref{trs41} and the system \eqref{trs43} reduces to, respectively,
\begin{equation}
\v=\lambda_0T+\cos \theta N+\lambda_1 B, \quad |\v |=1, \label{trs44}
\end{equation}
and
\begin{equation}
\left. 
\begin{array}{r}
f(1-\lambda_0^2)-\lambda'_0+ \kappa\cos\theta=0 \\ [4pt]
-f\lambda_0\cos\theta+\theta'\sin \theta-\lambda_0\kappa+\lambda_1\tau=0 \\ [4pt]
f\lambda_0\lambda_{1}+\lambda'_{1}+\cos\theta\tau=0.
\end{array}%
\right. \label{trs45}
\end{equation}

\begin{equation}
\v=\lambda_0T+\cos \theta N+\lambda_1 B, \quad |\v |=1, 
\end{equation}

\begin{theorem}\label{thm42}
Let $\gamma $ be a Frenet curve in a connected orientable Riemannian $3$-manifold and $\v$ a torse-forming vector field along $\gamma$ with potential function $f$. Then $\gamma$ is a PA curve with $(\v,\frac{\pi}{2})$ if and only if 
$$
\v=(\tanh p) T\pm (\sech p) B, \quad p:=\int^s f(u)du+r, \, r\in \r.
$$
In addition, the ratio of the curvatures of $\gamma$ is $\frac{\tau}{\kappa}=\sinh (p)$ and in particular this ratio is nonconstant.
\end{theorem}

\begin{proof}
If $\gamma$ is a PA curve with $(\v,\frac{\pi}{2})$, then $\langle\v,N\rangle=0$. By \eqref{trs44}, there is a smooth function $\vartheta$ defined along $\gamma$ such that $\v=\cos \vartheta T+\sin\vartheta B$.
Moreover, the system \eqref{trs45} is now
\begin{equation}
\left. 
\begin{array}{r}
\sin\vartheta(\vartheta'+f\sin\vartheta)=0 \\ [4pt]
\tau\sin\vartheta=\kappa \cos\vartheta \\ [4pt]
\cos\vartheta\left(\vartheta'+f\sin\vartheta\right)=0.
\end{array}%
\right. \label{simp41}
\end{equation}
If $\cos\vartheta=0$, then, by the first equation in \eqref{simp41}, $f$ would vanish along $\gamma$. Hence, from Lemma \ref{consleng} it follows that $\v$ becomes parallel along $\gamma$, which is not our case. Similarly, from the second equation in \eqref{simp41}, the case $\sin\vartheta=0$ is not possible because  $\gamma$ is assumed to be a Frenet curve. Consequently, it follows from \eqref{simp41} that
$$
\vartheta'+f\sin\vartheta=0.
$$
Integrating this equation, we derive $\cos\vartheta=\tanh p$ and $ \sin\vartheta=\sech p.$ Finally, the ratio of $\kappa$ and $\tau$ is concluded by the second equation in \eqref{simp41}, which is non-constant because $f$ is non-vanishing along $\gamma$. The proof of the converse follows from direct calculations and Theorem \ref{prop41}.

\end{proof}

\end{document}
